\section{De los modelos de lenguaje a los chatbots}

\subsection{Comportamiento de los modelos de lenguaje originales}

El profesor usa un modelo LaMDA más antiguo de Google para la demostración, porque todavía no ha pasado por mucho fine-tuning de post-training, lo que facilita observar el comportamiento original de next-word prediction.

Cuando se introduce directamente ``Hi, do you have any recommendations for dinner?'', la salida del modelo está lejos de parecer la de un chatbot competente. En efecto menciona algunas recomendaciones de restaurantes (como ``The Fat Duck''), pero después produce un texto sin ninguna relación como ``TripAdvisor staff removed this post''.

\begin{importantbox}{El modelo de lenguaje es una ``búsqueda difusa'' de los datos de entrenamiento}
En esencia, el modelo de lenguaje realiza una búsqueda difusa (fuzzy lookup) de sus datos de entrenamiento. Es muy probable que LaMDA genere ``TripAdvisor staff removed this post'' porque los datos de entrenamiento contienen contenido de los foros de TripAdvisor, y este texto estandarizado aparece con una frecuencia extremadamente alta en las publicaciones de esos foros, por lo que el modelo tiende a generarlo. Esto no es el modelo ``pensando'', sino reproduciendo los patrones estadísticos de los datos de entrenamiento.
\end{importantbox}

\subsection{Construcción gradual de una interfaz de chat}

El profesor mostró el proceso de transformación gradual, de un modelo de lenguaje original a un chatbot utilizable:

\textbf{Paso 1: Role Prompting (prompting de rol).} Se agrega ``You are a helpful chatbot'' antes de la entrada del usuario, para dirigir al modelo hacia la región de texto ``útil'' de los datos de entrenamiento. El efecto mejora un poco, pero el modelo interpreta a la vez a ambas partes de la conversación.

\textbf{Paso 2: prompting de formato.} Se usa un formato similar al de un guion de película, ``User: ...\textbackslash n Chatbot: ...'', para indicarle al modelo que distinga los roles de la conversación. El modelo capta de inmediato la indicación de formato, e incluso se pone un nombre a sí mismo (``HelpBot'').

\begin{knowledgebox}{¿Por qué funciona el prompting de formato?}
En los datos de entrenamiento del modelo existe una gran cantidad de datos de conversación con formato (publicaciones de foros, guiones de película, registros de atención al cliente, etcétera), y estos datos suelen distinguir a los hablantes con etiquetas del tipo ``User:'' / ``Assistant:''. Al usar un formato similar, dirigimos al modelo hacia la región del espacio de probabilidad que corresponde al texto con formato de conversación dentro de los datos de entrenamiento. El modelo no ``entiende'' el concepto de rol, simplemente genera el texto que con mayor probabilidad sigue a ese formato.
\end{knowledgebox}

\textbf{Paso 3: truncar la salida.} El modelo todavía genera la conversación de ambas partes a la vez. La solución es sencilla: conservar solo el texto entre ``Chatbot:'' y la siguiente etiqueta ``User:'', y descartar el resto.

\textbf{Paso 4: construir un harness de conversación.} Se escribe un programa de gestión de la conversación que mantiene el historial completo de la conversación; cada vez se concatena el historial al prompt que se envía al modelo, y al recibir la respuesta se agrega al historial, repitiendo el ciclo.

\begin{warningbox}{El modelo no está ``conversando''}
El profesor enfatiza especialmente que todo este proceso ``no es el preludio de Terminator 2: Rise of the Machines''. Que el modelo genere a la vez la conversación de ambas partes no se debe a que tenga conciencia de sí mismo, sino a que sus datos de entrenamiento contienen una gran cantidad de textos de conversación completos (guiones de película, etcétera); solo está haciendo next-word prediction. Necesitamos medios de ingeniería (prompting de formato + truncamiento + harness) para restringir su comportamiento.
\end{warningbox}

\subsection{Resumen de la sección}

Los pasos clave para pasar de un modelo de lenguaje original a un chatbot son: (1) dirigir el comportamiento del modelo mediante role prompting; (2) ayudar al modelo a distinguir los roles mediante prompting de formato; (3) lograr una conversación interactiva de varios turnos mediante el truncamiento y un harness de conversación. En esencia, todo el proceso consiste en usar el prompt para dirigir al modelo hacia la región del espacio de probabilidad de los datos de entrenamiento que deseamos.
